\documentclass[aspectratio=169,11pt]{beamer}
\usetheme{default}
\setbeamertemplate{navigation symbols}{}
\setbeamertemplate{footline}{%
  \hfill\insertframenumber/\inserttotalframenumber\hspace*{4pt}\vspace{2pt}%
}
\usepackage{lmodern}
\usepackage[most]{tcolorbox}
\usepackage{listings}
\usepackage{tikz}
\usetikzlibrary{arrows.meta,positioning,shapes.geometric,fit,backgrounds}
\usepackage{booktabs}
\usepackage{hyperref}

\definecolor{lcgreen}{HTML}{2E7D32}
\definecolor{lcblue}{HTML}{1565C0}
\definecolor{lcorange}{HTML}{EF6C00}
\definecolor{lcred}{HTML}{C62828}
\definecolor{codebg}{HTML}{F5F5F5}
\definecolor{docgray}{gray}{0.55}

\newtcolorbox{conceptbox}[1][]{colback=yellow!20,colframe=yellow!60!black,title=#1,fonttitle=\bfseries,top=2pt,bottom=2pt,left=4pt,right=4pt}
\newtcolorbox{warningbox}[1][]{colback=red!15,colframe=red!60!black,title=#1,fonttitle=\bfseries,top=2pt,bottom=2pt,left=4pt,right=4pt}
\newtcolorbox{tipbox}[1][]{colback=green!10,colframe=green!50!black,title=#1,fonttitle=\bfseries,top=2pt,bottom=2pt,left=4pt,right=4pt}
\newtcolorbox{codebox}[1][]{colback=blue!10,colframe=blue!40!black,title=#1,fonttitle=\bfseries,top=2pt,bottom=2pt,left=4pt,right=4pt}

\lstset{
  basicstyle=\ttfamily\scriptsize,
  backgroundcolor=\color{codebg},
  breaklines=true,
  frame=single,
  rulecolor=\color{blue!40!black},
  keywordstyle=\color{lcblue}\bfseries,
  stringstyle=\color{lcgreen},
  commentstyle=\color{docgray}\itshape,
  language=Python,
  showstringspaces=false,
  tabsize=2,
  xleftmargin=2pt,
  xrightmargin=2pt,
  aboveskip=2pt,
  belowskip=2pt,
}

\newcommand{\docref}[1]{{\scriptsize\color{docgray}[Docs: #1]}}

\title{LangGraph v1.0.x}
\subtitle{Section 02: Building Your First Graph}
\author{Agentic AI Foundations}
\date{March 2026}

\begin{document}

\begin{frame}
\titlepage
\end{frame}

\section{Mental Model}
\begin{frame}[t]{From Concept to Running Graph}
\setlength{\itemsep}{1pt}
\vspace{0.1cm}

\begin{conceptbox}[The Five-Step Pattern]
\setlength{\itemsep}{1pt}
\begin{enumerate}
\setlength{\itemsep}{1pt}
  \item \textbf{Define} state schema (TypedDict)
  \item \textbf{Create} StateGraph instance
  \item \textbf{Add} nodes (functions)
  \item \textbf{Add} edges (control flow)
  \item \textbf{Compile} and invoke
\end{enumerate}
\end{conceptbox}

\vspace{0.15cm}

This same five-step pattern scales from ``hello world'' to production agents.
\end{frame}

\section{Step-by-Step}
\begin{frame}[fragile,t]{Step 1: Define State Schema}
\vspace{0.1cm}

\begin{codebox}[TypedDict]
\begin{lstlisting}
from typing_extensions import TypedDict

class State(TypedDict):
    user_input: str
    response: str
\end{lstlisting}
\end{codebox}

\vspace{0.1cm}

\begin{tipbox}[State Design]
Think about what data flows through your graph. Each field is a piece of state.
\end{tipbox}
\end{frame}

\begin{frame}[fragile,t]{Step 2: Create StateGraph}
\vspace{0.1cm}

\begin{codebox}[Initialize]
\begin{lstlisting}
from langgraph.graph import StateGraph

graph = StateGraph(State)
\end{lstlisting}
\end{codebox}

\vspace{0.15cm}

That's it. The StateGraph holds your nodes and edges.

\vspace{0.15cm}

\begin{tipbox}[Methods]
\texttt{graph.add\_node()}, \texttt{graph.add\_edge()}, \texttt{graph.compile()}
\end{tipbox}
\end{frame}

\begin{frame}[fragile,t]{Step 3 \& 4: Add Nodes and Edges}
\vspace{0.05cm}

\begin{codebox}[Define a Node]
\begin{lstlisting}
def process_input(state: State) -> dict:
    response = state["user_input"].upper()
    return {"response": response}

graph.add_node("process", process_input)
\end{lstlisting}
\end{codebox}

\vspace{0.1cm}

\begin{codebox}[Connect with Edges]
\begin{lstlisting}
from langgraph.graph import START, END
graph.add_edge(START, "process")
graph.add_edge("process", END)
\end{lstlisting}
\end{codebox}
\end{frame}

\begin{frame}[fragile,t]{Step 5: Compile and Invoke}
\vspace{0.1cm}

\begin{codebox}[Compile]
\begin{lstlisting}
compiled_graph = graph.compile()
\end{lstlisting}
\end{codebox}

\vspace{0.1cm}

\begin{codebox}[Invoke]
\begin{lstlisting}
result = compiled_graph.invoke(
    {"user_input": "hello"}
)
print(result["response"])  # HELLO
\end{lstlisting}
\end{codebox}

\vspace{0.1cm}

{\footnotesize \texttt{invoke()} runs the graph synchronously and returns final state.}
\end{frame}

\section{StateGraph Class}
\begin{frame}[fragile,t]{StateGraph: Constructor \& Methods}
\vspace{0.05cm}

\begin{codebox}[Constructor]
\begin{lstlisting}
StateGraph(State)  # Takes your state schema
\end{lstlisting}
\end{codebox}

\vspace{0.1cm}

\begin{tabular}{p{4.5cm}p{6.5cm}}
\toprule
\textbf{Method} & \textbf{Purpose} \\
\midrule
\texttt{add\_node(name, fn)} & Add a computation node \\
\texttt{add\_edge(src, dst)} & Add a normal edge \\
\texttt{add\_conditional\_edges()} & Add conditional routing \\
\texttt{compile()} & Finalize and return runnable \\
\bottomrule
\end{tabular}

\vspace{0.15cm}

{\footnotesize Always call \texttt{compile()} before \texttt{invoke()}.}
\end{frame}

\section{Adding Nodes}
\begin{frame}[fragile,t]{Adding Nodes: graph.add\_node()}
\vspace{0.1cm}

\begin{conceptbox}[Node Registration]
\texttt{add\_node(name, function)} registers a named node.
\end{conceptbox}

\vspace{0.1cm}

\begin{codebox}[Multiple Nodes]
\begin{lstlisting}
graph.add_node("node_a", func_a)
graph.add_node("node_b", func_b)
graph.add_node("node_c", func_c)
\end{lstlisting}
\end{codebox}

\vspace{0.1cm}

\begin{tipbox}[Names]
Node names are strings. Use them in \texttt{add\_edge()} calls.
\end{tipbox}
\end{frame}

\section{Adding Edges}
\begin{frame}[fragile,t]{Adding Edges: graph.add\_edge()}
\vspace{0.1cm}

\begin{codebox}[Normal Edges]
\begin{lstlisting}
graph.add_edge(START, "node_a")
graph.add_edge("node_a", "node_b")
graph.add_edge("node_b", END)
\end{lstlisting}
\end{codebox}

\vspace{0.15cm}

\begin{tipbox}[Order Matters]
Data flows from source to destination. Define nodes before adding edges!
\end{tipbox}
\end{frame}

\section{Compilation}
\begin{frame}[fragile,t]{Compiling: graph.compile()}
\vspace{0.1cm}

\begin{conceptbox}[What compile() does]
\setlength{\itemsep}{1pt}
\begin{itemize}
  \item Validates graph structure (reachable, valid edges)
  \item Returns a \textit{runnable} (has \texttt{invoke()}, \texttt{stream()})
  \item Prepares execution engine
\end{itemize}
\end{conceptbox}

\vspace{0.15cm}

\begin{codebox}[Compile]
\begin{lstlisting}
compiled = graph.compile()
# compiled is now a Runnable
\end{lstlisting}
\end{codebox}

\vspace{0.15cm}

{\footnotesize Compile only once, invoke many times.}
\end{frame}

\section{Invoking}
\begin{frame}[fragile,t]{Invoking: compiled.invoke()}
\vspace{0.1cm}

\begin{conceptbox}[Execution]
\texttt{invoke(initial\_state)} runs the graph from START to END.
\end{conceptbox}

\vspace{0.1cm}

\begin{codebox}[Invoke and Return]
\begin{lstlisting}
result = compiled.invoke({
    "user_input": "hello world"
})
print(result)  # Full final state
\end{lstlisting}
\end{codebox}

\vspace{0.1cm}

\begin{tipbox}[Synchronous]
\texttt{invoke()} blocks until completion. For async, use \texttt{ainvoke()}.
\end{tipbox}
\end{frame}

\section{Complete Example: Chatbot}
\begin{frame}[fragile,t,shrink=3]{Complete Example: Simple Chatbot}
\vspace{0.05cm}

\begin{codebox}[Full Chatbot]
\begin{lstlisting}
from typing_extensions import TypedDict, Annotated
from langgraph.graph import StateGraph, START, END
from langgraph.graph import add_messages

class State(TypedDict):
    messages: Annotated[list, add_messages]

def chatbot_node(state: State):
    # Simulate LLM response
    return {"messages": ["Bot: Hello!"]}

graph = StateGraph(State)
graph.add_node("chatbot", chatbot_node)
graph.add_edge(START, "chatbot")
graph.add_edge("chatbot", END)
compiled = graph.compile()

result = compiled.invoke({
    "messages": ["User: Hi"]
})
\end{lstlisting}
\end{codebox}
\end{frame}

\section{Example 2: Classification}
\begin{frame}[fragile,t,shrink=4]{Example: Classification Pipeline}
\vspace{0.05cm}

\begin{codebox}[Classifier]
\begin{lstlisting}
class State(TypedDict):
    text: str
    category: str

def classify(state: State) -> dict:
    text = state["text"]
    cat = "spam" if "buy now" in text else "ok"
    return {"category": cat}

graph = StateGraph(State)
graph.add_node("classifier", classify)
graph.add_edge(START, "classifier")
graph.add_edge("classifier", END)
compiled = graph.compile()

result = compiled.invoke(
    {"text": "Check this out"}
)
\end{lstlisting}
\end{codebox}
\end{frame}

\section{Visualization}
\begin{frame}[fragile,t]{Getting Graph Visualization}
\vspace{0.1cm}

\begin{conceptbox}[Mermaid Diagram]
View your graph structure as a flowchart.
\end{conceptbox}

\vspace{0.1cm}

\begin{codebox}[Draw Mermaid]
\begin{lstlisting}
mermaid_str = compiled.get_graph().draw_mermaid()
print(mermaid_str)
# Paste into mermaid.live to visualize
\end{lstlisting}
\end{codebox}

\vspace{0.15cm}

\begin{tipbox}[Output]
ASCII flowchart showing all nodes and edges. Great for debugging!
\end{tipbox}
\end{frame}

\section{Config}
\begin{frame}[fragile,t]{Passing Config: RunnableConfig}
\vspace{0.1cm}

\begin{conceptbox}[Config dict]
Control graph execution behavior (threading, tags, streaming).
\end{conceptbox}

\vspace{0.1cm}

\begin{codebox}[Example]
\begin{lstlisting}
from langgraph.types import RunnableConfig

config = {"tags": ["production"]}
result = compiled.invoke(
    {"text": "test"},
    config=config
)
\end{lstlisting}
\end{codebox}

\vspace{0.15cm}

{\footnotesize Common config keys: \texttt{tags}, \texttt{run\_name}, \texttt{recursion\_limit}.}
\end{frame}

\section{Schemas}
\begin{frame}[fragile,t,shrink=3]{Input and Output Schemas}
\vspace{0.1cm}

\begin{tipbox}[Automatic]
LangGraph infers input/output schema from your State TypedDict.
\end{tipbox}

\vspace{0.1cm}

\begin{codebox}[Inspect Schemas]
\begin{lstlisting}
print(compiled.input_schema)
print(compiled.output_schema)
\end{lstlisting}
\end{codebox}

\vspace{0.15cm}

\begin{conceptbox}[Why It Matters]
Used for validation, API documentation, and tool-calling in frameworks.
\end{conceptbox}
\end{frame}

\section{Type Checking}
\begin{frame}[fragile,t]{Type Checking and Validation}
\vspace{0.1cm}

\begin{conceptbox}[TypedDict Validation]
Use \texttt{typing\_extensions} for runtime validation.
\end{conceptbox}

\vspace{0.1cm}

\begin{codebox}[Validation]
\begin{lstlisting}
# Invoke with wrong type
result = compiled.invoke({
    "text": 123  # Should be string
})  # May error or coerce
\end{lstlisting}
\end{codebox}

\vspace{0.15cm}

\begin{warningbox}[Tip]
Always define state fields with correct types. Type hints prevent bugs!
\end{warningbox}
\end{frame}

\section{Anatomy of Graph}
\begin{frame}[t]{Anatomy of a Compiled Graph}
\vspace{0.1cm}

\begin{tikzpicture}[node distance=1.5cm, scale=0.9, every node/.style={font=\scriptsize}]
  \node[draw, fill=lcblue!30, rounded rectangle] (s) {START};
  \node[draw, fill=lcgreen!30, rounded rectangle, right=of s] (n1) {Node A};
  \node[draw, fill=lcgreen!30, rounded rectangle, right=of n1] (n2) {Node B};
  \node[draw, fill=lcred!30, rounded rectangle, right=of n2] (e) {END};

  \path[->, thick] (s) -- (n1);
  \path[->, thick] (n1) -- (n2);
  \path[->, thick] (n2) -- (e);

  \node[below=of s, font=\tiny] {Entry};
  \node[below=of n1, font=\tiny] {Compute};
  \node[below=of n2, font=\tiny] {Compute};
  \node[below=of e, font=\tiny] {Exit};
\end{tikzpicture}

\vspace{0.15cm}

Every compiled graph has this structure: START -> nodes → END.
\end{frame}

\section{Debug Tips}
\begin{frame}[fragile,t]{Debug Tips: Printing State}
\vspace{0.1cm}

\begin{codebox}[Debug Node]
\begin{lstlisting}
def debug_node(state: State) -> dict:
    print(f"State: {state}")
    # Do work
    return {"field": value}

graph.add_node("debug", debug_node)
\end{lstlisting}
\end{codebox}

\vspace{0.15cm}

\begin{tipbox}[Using stream()]
For better debugging, use \texttt{stream()} to see intermediate outputs.
\end{tipbox}
\end{frame}

\section{Common Patterns}
\begin{frame}[t]{Common Graph Patterns}
\vspace{0.1cm}

\begin{tipbox}[Linear Graph]
START -> A → B → C → END. Sequential processing.
\end{tipbox}

\vspace{0.1cm}

\begin{tipbox}[Branching Graph]
START -> A → (B or C or D) → END. Uses conditional edges.
\end{tipbox}

\vspace{0.1cm}

\begin{tipbox}[Looping Graph]
A -> (A again or next). Cycles for agent loops.
\end{tipbox}

\vspace{0.1cm}

{\footnotesize Section 03 covers branching and looping in detail.}
\end{frame}

\section{Summary}
\begin{frame}[t,shrink=4]{Summary: Building Your First Graph}
\setlength{\itemsep}{1pt}
\vspace{0.05cm}

\begin{conceptbox}[Five-Step Pattern]
\setlength{\itemsep}{0.5pt}
\begin{enumerate}
\setlength{\itemsep}{0.5pt}
  \item Define state (TypedDict)
  \item Create StateGraph(State)
  \item Add nodes with add\_node()
  \item Add edges with add\_edge()
  \item Compile and invoke()
\end{enumerate}
\end{conceptbox}

\vspace{0.1cm}

\begin{tipbox}[Key Methods]
\texttt{add\_node}, \texttt{add\_edge}, \texttt{compile}, \texttt{invoke}
\end{tipbox}

\vspace{0.1cm}

\begin{tipbox}[Debugging}
Use mermaid diagram, print state, test with simple inputs first.
\end{tipbox}
\end{frame}

\begin{frame}[t,shrink=5]{Self-Check Questions}
\setlength{\itemsep}{0pt}
\vspace{0.05cm}

{\footnotesize
\begin{enumerate}
\setlength{\itemsep}{0.5pt}
  \item What are the five steps to build a LangGraph?
  \item What does StateGraph constructor take?
  \item How do you add a node?
  \item When must you call compile()?
  \item What does invoke() return?
  \item How can you visualize your graph?
  \item What's the difference between invoke() and stream()?
\end{enumerate}

\vspace{0.05cm}

\textbf{Quick Answers:} (1) Define state, create graph, add nodes, add edges, compile/invoke. (2) State schema. (3) graph.add\_node(name, fn). (4) Before invoke(). (5) Final state dict. (6) get\_graph().draw\_mermaid(). (7) invoke blocks; stream returns events.
}
\end{frame}

\end{document}
